\section{Bases}

\setcounter{thm}{25}
\begin{mydef} [basis]
  A \qt{basis} $\beta$ of $V$ is a list of vectors $\beta = \onetilln{v}$ in $V$ that is linearly independent and spans $V$.
  \begin{equation}
    V = \myspan{v_1, \dots, v_n} = \myspan{ \beta }
  \end{equation}
\end{mydef}

\setcounter{thm}{27}
\begin{thm} [condition for a basis]
  A list $v_1, \dots, v_n$ of vectors in $V$ is a basis of $V$ $\iff$ every $v \in V$ can be written uniquely in the form
  \begin{equation}
    v = a_1 v_1 + \dots + a_n v_n \quad \text{for some scalars } a_1, \dots, a_n \in \myF.
  \end{equation}

  Note that the actual definition of linear independence says nothing about uniqueness.
\end{thm}
\begin{prf}
  \begin{description}
    \item{\Rightarrowdirection} Suppose $v_1, \dots, v_n$ is a basis of $V$. Since $v_1, \dots, v_n$ spans $V$, every $v \in V$ can be written in the form $v = a_1 v_1 + \dots + a_n v_n$ for some $a_1, \dots, a_n \in \myF$. To show uniqueness, suppose $v = c_1 v_1 + \dots + c_n v_n$ is another representation. Subtracting gives
    \begin{equation}
      0 = (a_1 - c_1) v_1 + \dots + (a_n - c_n) v_n.
    \end{equation}
    Since $v_1, \dots, v_n$ is linearly independent, $a_1 - c_1 = 0, \dots, a_n - c_n = 0$, so $a_1 = c_1, \dots, a_n = c_n$.

    \item{\Leftarrowdirection} Suppose every $v \in V$ can be written uniquely as a linear combination of $v_1, \dots, v_n$. Then $v_1, \dots, v_n$ spans $V$. Furthermore, the unique representation of $0 \in V$ must be $0 = 0v_1 + \dots + 0v_n$, which implies that $v_1, \dots, v_n$ is linearly independent. Thus, $v_1, \dots, v_n$ is a basis of $V$.
  \end{description}
  \vspace{-\baselineskip}
\end{prf}

\setcounter{thm}{29}
\begin{thm} [every spanning list contains a basis]
  \label{thm: every spanning list contains a basis}
  Every spanning list $v_1, \ddd, v_n$ of a vector space $V$ can be reduced to a basis of $V$. So if $V= \myspan {v_1, \ddd, v_n}$, then we can delete elements from that list, or leave the list unchanged, such that the list forms a basis. That is, some sublist of $v_1, \dots, v_n$ is a basis of $V$.
\end{thm}
\begin{prf}
  Suppose $V = \myspan{v_1, \ldots, v_n}$. Let $B_0 = (v_1, \ldots, v_n)$. We proceed in $n$ steps:

  \emph{\bfseries Step 1: }
  If $v_1 = 0$, delete $v_1$ from $B_0$:
  \begin{equation}
    \text{Therefore, let } B_1 :\equiv (v_2, \ldots, v_n).
  \end{equation}

  If $v_1 \neq 0$, then leave $B_0$ unchanged:
  \begin{equation}
    B_1 :\equiv B_0=(v_1, \ldots, v_n).
  \end{equation}

  \emph{\bfseries Step k: }
  Let $B_{k-1} = (\underbrace{v_{1}^*, \ldots, v_{k-1}^*}_{\substack{\text{modified} \\ \text{elements}}}, \underbrace{v^{\vphantom{*}}_{k}, \ldots, v^{\vphantom{*}}_n}_{\substack{\text{original} \\ \text{elements}}})$ be our list from our previous steps, where $v_1^*, \ldots, v_{k-1}^*$ are not necessarily the original elements $v_1, \ldots, v_{k-1}$ from $B_0$. The $*$ indicates that those elements may have been modified in the previous steps.

  If $v_k \in \myspan{v_1^*, \ldots, v_{k-1}^*}$, then delete $v_k$ from list $B_{k-1}$ such that:
  \begin{equation}
    B_k :\equiv (v_{1}^*, \ldots, v_{k-1}^*, v_{k+1}, \ldots, v_n).
  \end{equation}

  If $v_k \notin \myspan{v_1^*, \ldots, v_{k-1}^*}$, leave the list unchanged:
  \begin{equation}
    B_k :\equiv B_{k-1}.
  \end{equation}

  Stop the process after step $n$, getting a list $B_n$. By the Linear Dependence Lemma (\ref{thm: linear dependence lemma}), deleting a vector that is in the span of the previous vectors does not change the span, so $B_n$ still spans $V$. Furthermore, the construction guarantees that no vector in $B_n$ is in the span of the previous ones. By the Linear Dependence Lemma (\ref{thm: linear dependence lemma}), this implies that $B_n$ is linearly independent (otherwise some vector would be in the span of its predecessors). Hence, $B_n$ is a basis of $V$.
  %TODO: Why does axler cite linear dependence lemma 2.19.
\end{prf}

\setcounter{thm}{30}
\begin{thm} [basis of finite\-/dimensional vector space]
  \label{thm: every finite-dimensional vector space has a basis}
  Every \findimvs has a basis.
\end{thm}
\begin{prf}
  By definition, a \fdvs has a spanning list. By \ref{thm: every spanning list contains a basis}, this spanning list can be reduced to a basis of $V$.
\end{prf}

\begin{thm}
  \label{thm: every linearly independent list of vectors in a finite-dimensional vector space can be extended to a basis of the vector space}
  Every \lid list of vectors in a \findimvs can be extended to a basis of the \vs.
\end{thm}
\begin{prf}
  Let $u_1, \ldots, u_m \in V$ be linearly independent, and let $w_1, \ddd, w_n$ be a spanning list of $V$. Then the joined list
  \begin{equation}
    u_1, \ddd , u_m, w_1, \ddd , w_n
  \end{equation}
  spans $V$. We apply the reduction procedure of \ref{thm: every spanning list contains a basis} to this list.
  Since $u_1, \dots, u_m$ is linearly independent, no $u_j$ is in the span of the preceding vectors $u_1, \dots, u_{j-1}$. Thus, none of the $u_j$'s are deleted during the first $m$ steps.
  The algorithm will delete only some of the $w_k$'s, resulting in a basis of $V$ that contains all the original vectors $u_1, \dots, u_m$.
\end{prf}

\begin{thm} [every subspace of $V$ is part of a direct sum equal to $V$]
\label{thm:every-subspace-of-V-is-part-of-a-direct-sum-equal-to-V}
  Suppose $V$ is finite-dimensional and $U$ is a subspace of $V$. Then there is a subspace $W$ of $V$ such that
  \begin{equation}
    V = U \oplus W.
  \end{equation}
\end{thm}
\begin{prf}
  Suppose $U$ is a subspace of a finite-dimensional vector space $V$.
  By \ref{thm: finite-dimensional subspace}, $U$ is also finite-dimensional, so by \ref{thm: every finite-dimensional vector space has a basis}, $U$ has a basis $u_1, \ldots, u_m$.
  Since $u_1, \dots, u_m$ is linearly independent in $V$, by \ref{thm: every linearly independent list of vectors in a finite-dimensional vector space can be extended to a basis of the vector space}, it can be extended to a basis $u_1, \ldots, u_m, w_1, \ldots, w_n$ of $V$.
  Set
  \begin{equation}
    W := \myspan{w_1, \ldots, w_n}.
  \end{equation}

  To prove that $V = U \oplus W$, by \ref{thm: sum and intersection of two subspaces}, it suffices to show that $V = U + W$ and $U \cap W = \{0\}$.
  \begin{description}
    \item{\bfemph{\slshape Proof of $V = U + W$:}}
    Let $v \in V$. Since $u_1, \ldots, u_m, w_1, \ldots, w_n$ is a basis of $V$, there exist scalars $a_1, \dots, a_m, b_1, \dots, b_n \in \myF$ such that
    \begin{equation}
      v = \underbrace{(a_1 u_1 + \cdots + a_m u_m)}_{\in \, U} + \underbrace{(b_1 w_1 + \cdots + b_n w_n)}_{\in \, W}.
    \end{equation}
    Hence $v \in U + W$, showing $V \subseteq U + W$. Since $U, W \subseteq V$, we have $U + W \subseteq V$, so $V = U + W$.

    \item{\bfemph{\slshape Proof of $U \cap W = \{0\}$:}}
    Let $v \in U \cap W$. Since $v \in U$, $v = a_1 u_1 + \cdots + a_m u_m$ for some $a_j \in \myF$. Since $v \in W$, $v = b_1 w_1 + \cdots + b_n w_n$ for some $b_k \in \myF$.
    Subtracting gives
    \begin{equation}
      a_1 u_1 + \cdots + a_m u_m - b_1 w_1 - \cdots - b_n w_n = 0.
    \end{equation}
    Since $u_1, \dots, u_m, w_1, \dots, w_n$ is a basis of $V$, it is linearly independent, so all $a_j = 0$ and all $b_k = 0$. Thus $v = 0$, proving $U \cap W = \{0\}$.
  \end{description}
  Therefore, $V = U \oplus W$, as desired.
\end{prf}
